\documentclass[11pt,a4paper]{article}
\usepackage[T1]{fontenc}
\usepackage[utf8]{inputenc}
\usepackage{amsmath,amssymb,xcolor}
\usepackage[margin=25mm]{geometry}
\usepackage[colorlinks=true,urlcolor=blue]{hyperref}
\definecolor{nexorablue}{HTML}{2459B5}
\setlength{\parindent}{0pt}
\setlength{\parskip}{7pt}
\setlength{\emergencystretch}{2em}
\renewcommand{\familydefault}{\sfdefault}
\begin{document}

{\small\bfseries\color{nexorablue} NEXORA CPGE}\par
{\small MPSI / NEX-RGT-06 / Énoncé}\par
{\LARGE\bfseries Bilan chiffré d’une charge complète}\par
\vspace{0.25em}\hrule\vspace{0.6em}

Une capacité initialement vide, $C=220\,\mathrm{\mu F}$, est chargée sous $E=9{,}0\,\mathrm V$ par une résistance positive. Calculer l’énergie fournie, l’énergie stockée à la fin et l’énergie dissipée. Expliquer ce qui varie si l’on double $R$.


\end{document}
